\begin{theorem}[P.\ Gabriel]\label{prop:identify-ModR}
	设 $s \in \Obj(\mathcal{A})$, 命 $R := \End_{\mathcal{A}}(s)$, 则
	\[ \Hom_{\mathcal{A}}(s, \cdot): \mathcal{A} \to \cated{Mod}R \]
	为等价的充要条件是 $s$ 为 $\mathcal{A}$ 的紧投射生成元.
\end{theorem}
\begin{proof}
	简记 $G := \Hom_{\mathcal{A}}(s, \cdot): \mathcal{A} \to \cated{Mod}R$. 先说明条件的必要性. 由于 $G$ 是等价并且 $G(s) = R$, 问题归结为验证 $R$ 是 $\cated{Mod}R$ 的紧投射生成元, 这是简单的.
	
	以下处理充分性. 设 $s$ 是紧投射生成元. 命题 \ref{prop:proj-compact-gen} 蕴涵 $G$ 是保所有小 $\varinjlim$ 的忠实正合函子, 因为它与忘却 $\cated{Mod}R \to \Bbbk\dcate{Mod}$ 的合成有此性质. 既然忠实性已知, 问题归结为证:
	\begin{compactitem}
		\item 对所有 $X, Y \in \Obj(\mathcal{A})$, 相应的 $\Hom_{\mathcal{A}}(X, Y) \to \Hom_R(GX, GY)$ 是满射;
		\item $G$ 是本质满的.
	\end{compactitem}
	
	我们首先说明对于任何 $X \in \Obj(\mathcal{A})$, 定义态射 $\phi_X: s^{\oplus \Hom_{\mathcal{A}}(s, X)} \to X$ 使得它在 $f: s \to X$ 对应的直和项上是 $f$, 则 $\phi_X$ 是满的. 使用反证法: 若 $\Coker(\phi_X) \neq 0$, 则生成元的性质说明存在非零态射 $s \to \Coker(\phi_X)$ (命题 \ref{prop:injective-cogenerator}), 而投射对象的性质说明此态射可以分解为 $s \xrightarrow{g} X \to \Coker(\phi_X)$, 故 $s^{\oplus \Hom_{\mathcal{A}}(s, X)} \xrightarrow{\phi_X} X \to \Coker(\phi_X)$ 在 $g$ 对应的直和项上非零, 矛盾.
	
	对 $\Ker(\phi_X)$ 的核也可以继续如上操作, 由此得知任何 $X$ 皆可置入正合列
	\[ s^{\oplus J} \to s^{\oplus I} \to X \to 0, \quad I = I_X, \; J = J_X: \text{小集}. \]
	因为 $G$ 保所有小 $\varinjlim$, 这给出右 $R$-模的正合列
	\[ R^{\oplus J} \to R^{\oplus I} \to GX \to 0, \]
	两者的第一段态射都可以视同 $R$ 上的 $I \times J$ 矩阵, 以左乘作用; 它们实际是同一个矩阵.
	
	现在考虑右 $R$-模同态 $f: GX \to GY$. 对 $X$ 和 $Y$ 任取如上形式的正合列, 易见 $f$ 可提升为行正合交换图表
	\[\begin{tikzcd}
		R^{\oplus J_X} \arrow[r] \arrow[d] & R^{\oplus I_X} \arrow[d] \arrow[r] & GX \arrow[d, "f"] \arrow[r] & 0 \\
		R^{\oplus J_Y} \arrow[r] & R^{\oplus I_Y} \arrow[r] & GY \arrow[r] & 0.
	\end{tikzcd}\]
	但左侧方块的四个箭头都可以表成 $R$ 上的矩阵, 四个矩阵遂给出行正合交换图表
	\[\begin{tikzcd}
		s^{\oplus J_X} \arrow[r] \arrow[d] & s^{\oplus I_X} \arrow[d] \arrow[r] & X \arrow[dashed, d] \arrow[r] & 0 \\
		s^{\oplus J_Y} \arrow[r] & s^{\oplus I_Y} \arrow[r] & Y \arrow[r] & 0
	\end{tikzcd}\]
	的实线部分, 虚线部分由余核的函子性唯一确定. 由于 $G$ 正合, 必有 $G(X \dashrightarrow Y) = f$.
	
	类似的技巧可以说明 $G$ 本质满. 给定右 $R$-模 $M$, 存在正合列
	\[ R^{\oplus L} \to R^{\oplus K} \to M \to 0, \]
	其第一段映射视同 $R$ 上的 $K \times L$ 矩阵. 但同一个矩阵在 $\mathcal{A}$ 中给出余核正合列
	\[ s^{\oplus L} \to s^{\oplus K} \to X \to 0; \]
	既然 $G$ 保所有小 $\varinjlim$, 必有 $GX \simeq M$. 明所欲证.
\end{proof}
